\ifdefined\gkver\else\def\gkver{student}\fi
\documentclass[\gkver]{gaokaozhenti}
\begin{document}

\chapter{2005年北京理综}
%% paperType: 理综物理部分

\begin{enumerate}
\item
%% number: 14
%% paperName: 2005年北京理综第 14 题 6 分
%% typeId: 1
%% type: 单选题
%% chapter: 气体性质
%% point: 热力学第一定律
%% method: 无
%% score: 6
%% degree: 650
%% duplicateId: 0
%% body:
下列关于热现象的说法，正确的是\xzanswer{D}
\fourchoices[answer=D]
{外界对物体做功，物体的内能一定增加}
{气体的温度升高，气体的压强一定增大}
{任何条件下，热量都不会由低温物体传递到高温物体}
{任何热机都不可能使燃料释放的热量完全转化为机械能}
%% answer: D
%% memo:
\memoanswer{
本题原卷及材料中未提供详解，待补充。
}

\item
%% number: 15
%% paperName: 2005年北京理综第 15 题 6 分
%% typeId: 1
%% type: 单选题
%% chapter: 光学
%% point: 光的衍射
%% method: 无
%% score: 6
%% degree: 650
%% duplicateId: 0
%% body:
在下列各组的两个现象中都表现出光具有波动性的是\xzanswer{C}
\fourchoices[answer=C]
{光的折射现象、色散现象}
{光的反射现象、干涉现象}
{光的衍射现象、偏振现象}
{光的直线传播现象、光电效应现象}
%% answer: C
%% memo:
\memoanswer{
本题原卷及材料中未提供详解，待补充。
}

\item
%% number: 16
%% paperName: 2005年北京理综第 16 题 6 分
%% typeId: 1
%% type: 单选题
%% chapter: 原子和原子核
%% point: 质能方程
%% method: 无
%% score: 6
%% degree: 650
%% duplicateId: 0
%% body:
为纪念爱因斯坦对物理学的巨大贡献，联合国将 2005 年定为“国际物理年”。对于爱因斯坦提出的质能方程 $E=mc^{2}$，下列说法中不正确的是\xzanswer{D}
\fourchoices[answer=D]
{$E=mc^{2}$ 表明物体具有的能量与其质量成正比}
{根据 $\Delta E=\Delta mc^{2}$ 可计算核反应的能量}
{一个质子和一个中子结合成一个氘核时释放能量，表明此过程出现了质量亏损}
{$E=mc^{2}$ 中的 $E$ 是发生核反应中释放的核能}
%% answer: D
%% memo:
\memoanswer{
本题原卷及材料中未提供详解，待补充。
}

\item
%% number: 17
%% paperName: 2005年北京理综第 17 题 6 分
%% typeId: 1
%% type: 单选题
%% chapter: 机械波
%% point: 波的图象
%% method: 无
%% score: 6
%% degree: 650
%% duplicateId: 0
%% body:
一列简谐机械横波某时刻的波形图如图所示，波源的平衡位置坐标为 $x=0$。当波源质点处于其平衡位置上方且向下运动时，介质中平衡位置坐标 $x=2\Um$ 的质点所处位置及运动情况是\xzanswer{A}
\onepicture[w=8cm]{figs/17.png}
\fourchoices[answer=A]
{在其平衡位置下方且向上运动}
{在其平衡位置下方且向下运动}
{在其平衡位置上方且向上运动}
{在其平衡位置上方且向下运动}
%% answer: A
%% memo:
\memoanswer{
本题原卷及材料中未提供详解，待补充。
}

\item
%% number: 18
%% paperName: 2005年北京理综第 18 题 6 分
%% typeId: 2
%% type: 多选题
%% chapter: 交流电
%% point: 交流电
%% method: 无
%% score: 6
%% degree: 650
%% duplicateId: 0
%% body:
正弦交流电源与电阻 $R$、交流电压表按图 1 所示的方式连接，$R=10\UO$，交流表的示数是 $10\UV$。图 2 是交变电源输出电压 $u$ 随时间 $t$ 变化的图象，则\xzanswer{A}
\onepicture[w=12cm]{figs/18.png}
\fourchoices[answer=A]
{通过 $R$ 的电流 $i_{R}$ 随时间 $t$ 变化的规律是 $i_{R}=\sqrt{2}\cos100\pi t$（A）}
{通过 $R$ 的电流 $i_{R}$ 随时间 $t$ 变化的规律是 $i_{R}=\sqrt{2}\cos50\pi t$（A）}
{$R$ 两端的电压 $u_{R}$ 随时间 $t$ 变化的规律是 $u_{R}=5\sqrt{2}\cos100\pi t$（V）}
{$R$ 两端的电压 $u_{R}$ 随时间 $t$ 变化的规律是 $u_{R}=5\sqrt{2}\cos50\pi t$（V）}
%% answer: A
%% memo:
\memoanswer{
本题原卷及材料中未提供详解，待补充。
}

\item
%% number: 19
%% paperName: 2005年北京理综第 19 题 6 分
%% typeId: 1
%% type: 单选题
%% chapter: 直线运动
%% point: 匀速直线运动
%% method: 无
%% score: 6
%% degree: 550
%% duplicateId: 0
%% body:
一人看到闪电 $12.3\Us$ 后又听到雷声。已知空气中的声速约为 $330\Ums\sim340\Ums$，光速为 $3\times10^{8}\Ums$，于是他用 12.3 除以 3 很快估算出闪电发生位置到他的距离为 $4.1\Ukm$。根据你所学的物理知识可以判断\xzanswer{B}
\fourchoices[answer=B]
{这种估算方法是错误的，不可采用}
{这种估算方法可以比较准确地估算出闪电发生位置与观察考间的距离}
{这种估算方法没有考虑光的传播时间，结果误差很大}
{即使声速增大 2 倍以上，本题的估算结果依然正确}
%% answer: B
%% memo:
\memoanswer{
本题原卷及材料中未提供详解，待补充。
}

\item
%% number: 20
%% paperName: 2005年北京理综第 20 题 6 分
%% typeId: 1
%% type: 单选题
%% chapter: 曲线运动
%% point: 人造地球卫星
%% method: 无
%% score: 6
%% degree: 450
%% duplicateId: 0
%% body:
已知地球质量大约是月球质量的 81 倍，地球半径大约是月球半径的 4 倍。不考虑地球、月球自转的影响，由以上数据可推算出\xzanswer{C}
\fourchoices[answer=C]
{地球的平均密度与月球的平均密度之比约为 9:8}
{地球表面重力加速度与月球表面重力加速度之比约为 9:4}
{靠近地球表面沿圆轨道运行的航天器的周期与靠近月球表面沿圆轨道运行的航天器的周期之比约为 8:9}
{靠近地球表面沿圆轨道运行的航天器线速度与靠近月球表面沿圆轨道运行的航天器线速度之比约为 81:4}
%% answer: C
%% memo:
\memoanswer{
本题原卷及材料中未提供详解，待补充。
}

\item
%% number: 21
%% paperName: 2005年北京理综第 21 题 6 分
%% typeId: 1
%% type: 单选题
%% chapter: 电磁感应
%% point: 验证电磁感应实验
%% method: 无
%% score: 6
%% degree: 650
%% duplicateId: 0
%% body:
现将电池组、滑线变阻器、带铁芯的线圈 A、线圈 B、电流计及开关如下图连接，在开关闭合、线圈 A 放在线圈 B 中的情况下，某同学发现当他将滑线变阻器的滑动端 P 向左加速滑动时，电流计指针向右偏转。由此可以判断\xzanswer{B}
\onepicture[w=8cm]{figs/21.png}
\fourchoices[answer=B]
{线圈 A 向上移动或滑动变阻器滑动端 P 向右加速滑动，都能引起电流计指针向左偏转}
{线圈 A 中铁芯向上拔出或断开开关，都能引起电流计指针向右偏转}
{滑动变阻器的滑动端 P 匀速向左或匀速向右滑动，都能使电流计指针静止在中央}
{因为线圈 A、线圈 B 的绕线方向未知，故无法判断电流计指针偏转的方向}
%% answer: B
%% memo:
\memoanswer{
本题原卷及材料中未提供详解，待补充。
}

\item
%% number: 22
%% paperName: 2005年北京理综第 22 题 18 分
%% typeId: 4
%% type: 实验题
%% chapter: 电路
%% point: 多用表的使用
%% method: 无
%% score: 18
%% degree: 450
%% duplicateId: 0
%% body:
“黑盒子”表面有 a、b、c 三个接线柱，盒内总共有两个电子元件，每两个接线柱之间只可能连接一个元件。为了探明盒内元件的种类及连接方式，某位同学用多用电表进行了如下探测：

第一步：用电压挡，对任意两接线柱正、反向测量，指针不发生偏转。

第二步：用电阻 $\times1000$ 挡，对任意两个接线柱正、反向测量，指针偏转情况如下图所示。

\onepicture[w=13cm]{figs/22a.png}

\begin{enumerate}
	\item
	第一步测量结果表明盒内 \tkanswer[3cm]{不存在电源}。
	\item
	图 2 示出了图（a）和图（b）中欧姆表指针所处的位置，其对应的阻值是 \tkanswer[1.5cm]{1200} $\UO$，图 3 示出了图（c）中欧姆表指针所处的位置，其对应的阻值是 \tkanswer[1.5cm]{500} $\UO$。

	\twopicture[labelprefix={}, numA={图2}, numB={图3}, w=6.5cm]{figs/22b.png}{figs/22c.png}
	\item
	请在图 4 的接线柱间，用电路图符号画出盒内的元件及连接情况。

	\drawpicanswer{\onepicture[w=4.5cm]{figs/22d.png}}{\onepicture[w=4.5cm]{figs/22c_an.png}}
	\item
	一个小灯泡与 $3\UV$ 电池组的连接情况如图 5 所示。如果把图 5 中 e、f 两端用导线直接相连，小灯泡仍可正常发光。欲将 e、f 两端分别与黑盒子上的两个接线柱相连，使小灯泡仍可发光。那么，e 端应连接到 \tkanswer[1.5cm]{C} 接线柱，f 端应连接到 \tkanswer[1.5cm]{a} 接线柱。

	\onepicture[w=8cm]{figs/22e.png}
\end{enumerate}
%% answer:
\jdanswer{
\begin{enumerate}
	\item
	盒内不存在电源。
	\item
	$1200\UO$，$500\UO$。
	\item
	盒内元件及连接情况如图 4 所示，a、c 之间接一只二极管（a 为负极、c 为正极），b、c 之间接一个电阻。
	\item
	e 端应连接到 C 接线柱，f 端应连接到 a 接线柱。
\end{enumerate}
}
%% memo:
\memoanswer{
本题原卷及材料中未提供详解，待补充。
}

\item
%% number: 23
%% paperName: 2005年北京理综第 23 题 16 分
%% typeId: 6
%% type: 计算题
%% chapter: 功和能
%% point: 机械能守恒定律
%% method: 无
%% score: 16
%% degree: 550
%% duplicateId: 0
%% body:
AB 是竖直平面内的四分之一圆弧轨道，在下端 B 与水平直轨道相切，如图所示。一小球自 A 点起由静止开始沿轨道下滑。已知圆轨道半径为 $R$，小球的质量为 $m$，不计各处摩擦。求：
\begin{enumerate}
	\item
	小球运动到 B 点时的动能；
	\item
	小球下滑到距水平轨道的高度为 $\frac{R}{2}$ 时的速度大小和方向；
	\item
	小球经过圆弧轨道的 B 点和水平轨道的 C 点时，所受轨道支持力 $N_{B}$、$N_{C}$ 各是多大？
\end{enumerate}
\onepicture[w=5cm, align=right]{figs/23.png}
%% answer:
\jdanswer{
\begin{enumerate}
	\item
	$E_{k}=mgR$。
	\item
	$v=\sqrt{gR}$，速度方向沿圆弧的切线向下，与竖直方向成 $30^{\circ}$ 角。
	\item
	$N_{B}=3mg$，$N_{C}=mg$。
\end{enumerate}
}
%% memo:
\memoanswer{
（1）根据机械能守恒，小球运动到 B 点时的动能 $E_{k}=mgR$。

（2）根据机械能守恒，$\Delta E_{k}=\Delta E_{p}$，即 $\frac{1}{2}mv^{2}=\frac{1}{2}mgR$，解得小球速度大小 $v=\sqrt{gR}$，速度方向沿圆弧的切线向下，与竖直方向成 $30^{\circ}$ 角。

（3）根据牛顿运动定律及机械能守恒，在 B 点有 $N_{B}-mg=\frac{mv_{B}^{2}}{R}$，$mgR=\frac{1}{2}mv_{B}^{2}$，解得 $N_{B}=3mg$；在 C 点 $N_{C}=mg$。
}

\item
%% number: 24
%% paperName: 2005年北京理综第 24 题 18 分
%% typeId: 6
%% type: 计算题
%% chapter: 电场
%% point: 带电粒子的直线运动
%% method: 无
%% score: 18
%% degree: 450
%% duplicateId: 0
%% body:
真空中存在空间范围足够大的、水平向右的匀强电场。在电场中，若将一个质量为 $m$、带正电的小球由静止释放，运动中小球速度与竖直方向夹角为 $37^{\circ}$（取 $\sin37^{\circ}=0.6$，$\cos37^{\circ}=0.8$）。现将该小球从电场中某点以初速度 $v_{0}$ 竖直向上抛出。求运动过程中
\begin{enumerate}
	\item
	小球受到的电场力的大小及方向；
	\item
	小球从抛出点至最高点的电势能变化量；
	\item
	小球的最小动量的大小及方向。
\end{enumerate}
%% answer:
\jdanswer{
\begin{enumerate}
	\item
	$F_{e}=\frac{3}{4}mg$，方向水平向右。
	\item
	电势能减少 $\frac{9}{32}mv_{0}^{2}$。
	\item
	$p_{\min}=\frac{3}{5}mv_{0}$，方向与电场方向成 $37^{\circ}$ 角斜向上。
\end{enumerate}
}
%% memo:
\memoanswer{
（1）根据题设条件，小球由静止释放后沿与竖直方向成 $37^{\circ}$ 角的方向做匀加速直线运动，电场力大小 $F_{e}=mg\tan37^{\circ}=\frac{3}{4}mg$，电场力的方向水平向右。

（2）小球沿竖直方向做匀减速运动，速度 $v_{y}=v_{0}-gt$；沿水平方向做初速度为 0 的匀加速运动，加速度 $a_{x}=\frac{F_{e}}{m}=\frac{3}{4}g$。小球上升到最高点的时间 $t=\frac{v_{0}}{g}$，此过程小球沿电场方向位移 $s_{x}=\frac{1}{2}a_{x}t^{2}=\frac{3v_{0}^{2}}{8g}$，电场力做功 $W=F_{e}s_{x}=\frac{9}{32}mv_{0}^{2}$，故小球上升到最高点的过程中电势能减少 $\frac{9}{32}mv_{0}^{2}$。

（3）水平速度 $v_{x}=a_{x}t$，竖直速度 $v_{y}=v_{0}-gt$，小球的速度 $v=\sqrt{v_{x}^{2}+v_{y}^{2}}$。由以上各式得 $25g^{2}t^{2}/16-2v_{0}gt+(v_{0}^{2}-v^{2})=0$，解得当 $t=\frac{16v_{0}}{25g}$ 时 $v$ 有最小值 $v_{\min}=\frac{3}{5}v_{0}$。此时 $v_{x}=\frac{12}{25}v_{0}$，$v_{y}=\frac{9}{25}v_{0}$，$\tan\theta=\frac{v_{y}}{v_{x}}=\frac{3}{4}$，即与电场方向夹角为 $37^{\circ}$ 斜向上。小球动量的最小值为 $p_{\min}=mv_{\min}=\frac{3}{5}mv_{0}$，最小动量的方向与电场方向夹角为 $37^{\circ}$，斜向上。
}

\item
%% number: 25
%% paperName: 2005年北京理综第 25 题 20 分
%% typeId: 6
%% type: 计算题
%% chapter: 运动和力
%% point: 动量守恒定律
%% method: 无
%% score: 20
%% degree: 350
%% duplicateId: 0
%% body:
下图是导轨式电磁炮实验装置示意图。两根平行长直金属导轨沿水平方向固定，其间安放金属滑块（即实验用弹丸）。滑块可沿导轨无摩擦滑行，且始终与导轨保持良好接触。电源提供的强大电流从一根导轨流入，经过滑块，再从另一导轨流回电源。滑块被导轨中的电流形成的磁场推动而发射。在发射过程中，该磁场在滑块所在位置始终可以简化为匀强磁场，方向垂直于纸面，其强度与电流的关系为 $B=kI$，比例常量 $k=2.5\times10^{-6}\UTA$。已知两导轨内侧间距 $l=1.5\Ucm$，滑块的质量 $m=30\Ug$，滑块沿导轨滑行 $5\Um$ 后获得的发射速度 $v=3.0\Ukms$（此过程视为匀加速运动）。
\begin{enumerate}
	\item
	求发射过程中电源提供的电流强度。
	\item
	若电源输出的能量有 $4\%$ 转换为滑块的动能，则发射过程中电源的输出功率和输出电压各是多大？
	\item
	若此滑块射出后随即以速度 $v$ 沿水平方向击中放在水平面上的砂箱，它嵌入砂箱的深度为 $s'$。设砂箱质量为 $M$，滑块质量为 $m$，不计砂箱与水平面之间的摩擦。求滑块对砂箱平均冲击力的表达式。
\end{enumerate}
\onepicture[w=10cm, align=right]{figs/25.png}
%% answer:
\jdanswer{
\begin{enumerate}
	\item
	$I=8.5\times10^{5}\UA$。
	\item
	$P=1.0\times10^{9}\UW$，$U=1.2\times10^{3}\UV$。
	\item
	$f=\frac{Mmv^{2}}{2s'(m+M)}$。
\end{enumerate}
}
%% memo:
\memoanswer{
（1）由匀加速运动公式 $a=\frac{v^{2}}{2s}=9\times10^{5}\Umsq$。由安培力公式和牛顿第二定律，有 $F=IBl=kI^{2}l$，$kI^{2}l=ma$，因此 $I=\sqrt{\frac{ma}{kl}}=8.5\times10^{5}\UA$。

（2）滑块获得的动能是电源输出能量的 $4\%$，即 $P\Delta t\times4\%=\frac{1}{2}mv^{2}$。发射过程中电源供电时间 $\Delta t=\frac{v}{a}=0.33\times10^{-2}\Us$，所需的电源输出功率 $P=\frac{mv^{2}/2}{\Delta t\times4\%}=1.0\times10^{9}\UW$。由功率 $P=IU$，解得输出电压 $U=\frac{P}{I}=1.2\times10^{3}\UV$。

（3）分别对砂箱和滑块用动能定理，有 $fs_{M}=\frac{1}{2}MV^{2}$，$fs_{m}=\frac{1}{2}mV^{2}-\frac{1}{2}mv^{2}$。由牛顿第三定律 $f=-f'$ 和相对运动 $s_{m}=s_{M}+s'$，由动量守恒 $mv=(m+M)V$，联立求得 $fs'=\frac{Mmv^{2}}{2(m+M)}$，故平均冲击力 $f=\frac{Mmv^{2}}{2s'(m+M)}$。
}

\end{enumerate}

\end{document}
